% =================================================================
%  CSE 815 -- Machine Learning
%  Module 1, Week 3, Part 1
%  Classification: Logistic Regression, Decision Boundaries,
%  the Logistic Cost Function and Gradient Descent
% =================================================================
\documentclass[aspectratio=169,10pt]{beamer}
\usepackage{cse815}

\title[CSE 815 -- Week 3, Part 1]{Classification with Logistic Regression}
\subtitle{The Sigmoid Model, Decision Boundaries, and a Cost Function that Works}
\author{Rokan Uddin Faruqui\\ \small Professor, Department of CSE}
\institute{\small University of Chittagong \\ \texttt{rufaruqui@cu.ac.bd}}
\date[Week 3 / Part 1]{Week 3 $\cdot$ Part 1}

\begin{document}

\begin{frame}[plain,noframenumbering]
  \titlepage
\end{frame}

% ---------------------------------------------------------------
\begin{frame}{From predicting a number to predicting a class}
  \begin{columns}[T,onlytextwidth]
    \begin{column}{0.5\textwidth}
      \begin{block}{Weeks 1--2 gave us}
        \[ f_{\vw,b}(\vx)=\vw\cdot\vx + b \in \R \]
        \[ J(\vw,b)=\frac{1}{2m}\sum_{i=1}^{m}\bigl(f_{\vw,b}(\xii{i})-\yii{i}\bigr)^2 \]
        with scaling, an $\alpha$ search and engineered features to make it run.
      \end{block}
      \footnotesize
      Every one of those predictions was a \emph{quantity}. Today the answer is a
      \emph{category}: benign or malignant, spam or not, fraud or legitimate.
    \end{column}
    \begin{column}{0.46\textwidth}
      \begin{keyidea}[Today]
        \footnotesize
        \begin{enumerate}
          \item Why linear regression breaks on classes  
          \item The sigmoid and the logistic model 
          \item Decision boundaries, linear and not 
          \item Why squared error fails here 
          \item Logistic loss and the simplified cost  
          \item Gradient descent, worked by hand 
        \end{enumerate}
      \end{keyidea}
      \footnotesize The recipe is unchanged --- model, cost, gradient descent. All three
      pieces get replaced.
    \end{column}
  \end{columns}
\end{frame}

% ===============================================================
\section{Motivations}

\begin{frame}{Problems whose answer is a category}
  \footnotesize
  \begin{center}
  \renewcommand{\arraystretch}{1.2}
  \begin{tabular}{@{}L{4.0cm}L{3.6cm}L{1.6cm}L{1.6cm}@{}}
    \toprule
    \textbf{Question} & \textbf{Input $\vx$} & \textbf{$y=0$} & \textbf{$y=1$}\\
    \midrule
    Is this e-mail spam?          & word counts, sender    & not spam & spam\\
    Is this transaction fraud?    & amount, time, location & genuine  & fraud\\
    Is this tumour malignant?     & size, age, cell shape  & benign   & malignant\\
    Will this student pass?       & attendance, marks      & no       & yes\\
    \bottomrule
  \end{tabular}
  \end{center}
  \vspace{1mm}
  \begin{columns}[T,onlytextwidth]
    \begin{column}{0.5\textwidth}
      \begin{defnbox}[Binary classification]
        \small $y$ takes exactly one of two values, written $0$ and $1$ by convention.
        \begin{itemize}
          \item $y=0$: the \textbf{negative class} --- absence of the thing.
          \item $y=1$: the \textbf{positive class} --- presence of it.
        \end{itemize}
      \end{defnbox}
    \end{column}
    \begin{column}{0.46\textwidth}
      \begin{pitfall}[``Positive'' is not ``good'']
        \small
        A malignant tumour is the positive class. The labels name what you are
        \emph{detecting}, not what you want. Choosing them the other way round is legal but
        inverts every coefficient's sign.
      \end{pitfall}
    \end{column}
  \end{columns}
\end{frame}

\begin{frame}{Why linear regression breaks}
  \begin{columns}[T,onlytextwidth]
    \begin{column}{0.5\textwidth}
      \centering {\scriptsize Fit a line, threshold at $0.5$ --- looks fine}\par
      \begin{tikzpicture}
        \begin{axis}[width=6.1cm,height=3.9cm,xlabel={tumour size},ylabel={$y$},
          xmin=0,xmax=11,ymin=-0.25,ymax=1.35,ytick={0,0.5,1},
          tick label style={font=\scriptsize},label style={font=\scriptsize}]
          \addplot[only marks,mark=*,mark size=1.8pt,cublue,forget plot]
            coordinates {(1,0)(2,0)(3,0)(4,0)(6,1)(7,1)(8,1)(9,1)};
          \addplot[curose,very thick,domain=0:11,forget plot]{0.1667*x-0.3333};
          \draw[dashed,cugray] (axis cs:0,0.5)--(axis cs:11,0.5);
          \draw[cuteal,very thick] (axis cs:5,-0.25)--(axis cs:5,1.35);
          \node[font=\tiny,cuteal,anchor=south west] at (axis cs:5.1,1.0) {boundary};
        \end{axis}
      \end{tikzpicture}
    \end{column}
    \begin{column}{0.48\textwidth}
      \centering {\scriptsize Add one large-but-malignant tumour --- it collapses}\par
      \begin{tikzpicture}
        \begin{axis}[width=6.1cm,height=3.9cm,xlabel={tumour size},ylabel={$y$},
          xmin=0,xmax=21,ymin=-0.25,ymax=1.35,ytick={0,0.5,1},
          tick label style={font=\scriptsize},label style={font=\scriptsize}]
          \addplot[only marks,mark=*,mark size=1.8pt,cublue,forget plot]
            coordinates {(1,0)(2,0)(3,0)(4,0)(6,1)(7,1)(8,1)(9,1)(20,1)};
          \addplot[curose,very thick,domain=0:21,forget plot]{0.0641*x+0.1282};
          \draw[dashed,cugray] (axis cs:0,0.5)--(axis cs:21,0.5);
          \draw[cuteal,very thick] (axis cs:5.8,-0.25)--(axis cs:5.8,1.35);
          \node[font=\tiny,cuteal,anchor=south west] at (axis cs:5.95,1.0) {boundary};
        \end{axis}
      \end{tikzpicture}
    \end{column}
  \end{columns}
  \vspace{1mm}
  \footnotesize
  The new point is not an error --- it is a perfectly ordinary large malignant tumour. Yet the
  least-squares line tilts to reach it, the $0.5$ crossing slides from $5.0$ to $5.8$, and the
  tumour at size $6$ --- correctly called malignant before --- is now called benign.
  \begin{keyidea}
    Squared error punishes a confident \emph{correct} prediction such as $f=2.4$ for $y=1$.
    Linear regression is also unbounded: $f(\vx)$ can return $-3$ or $+7$, which cannot be
    read as a class or as a probability.
  \end{keyidea}
\end{frame}

% ===============================================================
\section{The Logistic Regression Model}

\begin{frame}{The sigmoid function}
  \begin{columns}[T,onlytextwidth]
    \begin{column}{0.52\textwidth}
      \centering
      \begin{tikzpicture}
        \begin{axis}[width=6.6cm,height=4.5cm,xlabel={$z$},ylabel={$g(z)$},
          xmin=-6,xmax=6,ymin=-0.08,ymax=1.12,ytick={0,0.5,1},
          grid=major,grid style={culight},
          tick label style={font=\scriptsize},label style={font=\scriptsize}]
          \addplot[cublue,very thick,domain=-6:6,samples=120]{1/(1+exp(-x))};
          \draw[dashed,cugray] (axis cs:-6,1)--(axis cs:6,1);
          \draw[dashed,cugray] (axis cs:-6,0.5)--(axis cs:6,0.5);
          \addplot[only marks,mark=*,mark size=1.6pt,curose]coordinates{(0,0.5)};
        \end{axis}
      \end{tikzpicture}
    \end{column}
    \begin{column}{0.44\textwidth}
      \footnotesize
      \[ g(z) = \frac{1}{1+e^{-z}}, \qquad 0 < g(z) < 1 \]
      \begin{itemize}
        \item $g(0)=0.5$ exactly --- the curve's centre.
        \item $g(z)\to 1$ as $z\to\infty$, $g(z)\to 0$ as $z\to-\infty$; it never reaches
              either.
        \item $g(-z) = 1-g(z)$, so it is symmetric about $(0,0.5)$.
      \end{itemize}
      \vspace{1mm}
      \begin{center}\scriptsize
      \begin{tabular}{@{}rrrrrr@{}}
        \toprule
        $z$ & $-3$ & $-1$ & $0$ & $1$ & $3$\\
        $g(z)$ & 0.0474 & 0.2689 & 0.5000 & 0.7311 & 0.9526\\
        \bottomrule
      \end{tabular}
      \end{center}
      \vspace{1mm}
      Also called the \emph{logistic} function --- hence the name of the model.
    \end{column}
  \end{columns}
\end{frame}

\begin{frame}{The logistic regression model}
  \small
  Take the linear model of Week 2 and pass its output through the sigmoid:
  \[
    z = \vw\cdot\vx + b, \qquad
    f_{\vw,b}(\vx) \;=\; g(z) \;=\; \frac{1}{1+e^{-(\vw\cdot\vx+b)}} \;\in\;(0,1).
  \]
  \begin{columns}[T,onlytextwidth]
    \begin{column}{0.5\textwidth}
      \begin{defnbox}[What the output means]
        \footnotesize
        $f_{\vw,b}(\vx)$ is the model's estimate of
        \[ P(y=1 \mid \vx;\ \vw,b), \]
        the probability that $y=1$ for this input, given these parameters.
        Since $y$ is $0$ or $1$,
        \[ P(y=0\mid\vx) = 1 - f_{\vw,b}(\vx). \]
      \end{defnbox}
    \end{column}
    \begin{column}{0.46\textwidth}
      \footnotesize
      \textbf{Reading a prediction.} A patient with $f_{\vw,b}(\vx)=0.7$: the model says
      there is a $70\%$ chance the tumour is malignant, and therefore a $30\%$ chance it is
      benign.
      \vspace{1mm}
      \begin{keyidea}
        \footnotesize
        The model outputs a \emph{probability}, not a class. Turning it into a class needs one
        more decision --- the threshold --- and that decision is yours, not the algorithm's.
      \end{keyidea}
    \end{column}
  \end{columns}
  \vspace{1mm}
  \footnotesize
  Despite the name, logistic regression is a \textbf{classification} algorithm. The
  ``regression'' is historical: it regresses the log-odds $\log\frac{p}{1-p} = \vw\cdot\vx+b$,
  which is linear in the parameters.
\end{frame}

% ===============================================================
\section{Decision Boundaries}

\begin{frame}{From probability to prediction}
  \footnotesize
  Predict $\hat y = 1$ when $f_{\vw,b}(\vx)\ge 0.5$. Because $g(z)\ge 0.5$ exactly when
  $z\ge 0$, this is the same as
  \[
    \hat y = 1 \iff \vw\cdot\vx + b \ \ge\ 0 ,
    \qquad
    \hat y = 0 \iff \vw\cdot\vx + b \ <\ 0 .
  \]
  \begin{columns}[T,onlytextwidth]
    \begin{column}{0.48\textwidth}
      \begin{defnbox}[Decision boundary]
        \small The set of inputs where $\vw\cdot\vx+b=0$: the surface on which the model is
        exactly undecided. For a linear $z$ it is a straight line in $2$D, a plane in $3$D, a
        hyperplane in general.
      \end{defnbox}
    \end{column}
    \begin{column}{0.48\textwidth}
      \begin{keyidea}
        \small The boundary is a property of the \textbf{parameters}, not of the training
        data. Two data sets that yield the same $(\vw,b)$ have the same boundary; the data
        only determined which $(\vw,b)$ was learned.
      \end{keyidea}
    \end{column}
  \end{columns}
  \vspace{1mm}
  \begin{pitfall}
    $0.5$ is a \emph{default}, not a law. Screening for cancer you might predict malignant at
    $f\ge 0.2$, accepting false alarms to miss fewer tumours; approving a loan you might
    demand $f\ge 0.9$. Changing the threshold moves the boundary without retraining anything.
  \end{pitfall}
\end{frame}

\begin{frame}{A linear decision boundary, worked out}
  \begin{columns}[T,onlytextwidth]
    \begin{column}{0.46\textwidth}
      \centering
      \begin{tikzpicture}
        \begin{axis}[width=5.9cm,height=4.4cm,xlabel={$x_1$},ylabel={$x_2$},
          xmin=0,xmax=4,ymin=0,ymax=4,grid=major,grid style={culight},
          tick label style={font=\scriptsize},label style={font=\scriptsize}]
          \addplot[draw=none,fill=cuteal!12,forget plot]
            coordinates {(0,3)(3,0)(4,0)(4,4)(0,4)};
          \addplot[only marks,mark=o,mark size=2.2pt,cublue,thick,forget plot]
            coordinates {(0.5,0.7)(1.0,1.2)(1.5,0.6)(0.7,1.8)(2.0,0.5)};
          \addplot[only marks,mark=x,mark size=2.6pt,curose,thick,forget plot]
            coordinates {(2.4,2.2)(3.2,1.6)(2.0,2.9)(3.4,2.8)(2.8,1.4)};
          \addplot[cublue,very thick,domain=0:3]{3-x};
          \node[font=\tiny,cublue,anchor=west] at (axis cs:0.15,3.4) {$x_1+x_2=3$};
        \end{axis}
      \end{tikzpicture}\par
      {\scriptsize $\circ$~$y=0$\qquad $\times$~$y=1$}
    \end{column}
    \begin{column}{0.5\textwidth}
      \footnotesize
      Take $\vw=[1,1]$ and $b=-3$, so $z = x_1+x_2-3$.
      \[ \text{boundary:}\quad x_1+x_2-3=0 \;\Longleftrightarrow\; x_1+x_2=3 \]
      \begin{center}\scriptsize
      \renewcommand{\arraystretch}{1.15}
      \begin{tabular}{@{}ccrl@{}}
        \toprule
        $\vx$ & $z$ & $f_{\vw,b}(\vx)$ & prediction\\
        \midrule
        $(1,1)$ & $-1$ & $0.2689$ & $\hat y = 0$\\
        $(2,1)$ & $0$   & $0.5000$ & on the boundary\\
        $(2,2)$ & $+1$  & $0.7311$ & $\hat y = 1$\\
        $(3,3)$ & $+3$  & $0.9526$ & $\hat y = 1$, confidently\\
        \bottomrule
      \end{tabular}
      \end{center}
      \vspace{1mm}
      Distance from the boundary is confidence: $|z|$ large means $f$ near $0$ or $1$. The
      magnitude of $\vw$ controls how fast confidence grows as you step away.
    \end{column}
  \end{columns}
\end{frame}

\begin{frame}{Non-linear decision boundaries}
  \begin{columns}[T,onlytextwidth]
    \begin{column}{0.46\textwidth}
      \centering
      \begin{tikzpicture}
        \begin{axis}[axis equal image,width=5.0cm,xlabel={$x_1$},ylabel={$x_2$},
          xmin=-2,xmax=2,ymin=-2,ymax=2,grid=major,grid style={culight},
          tick label style={font=\scriptsize},label style={font=\scriptsize}]
          \addplot[only marks,mark=o,mark size=2.2pt,cublue,thick,forget plot]
            coordinates {(0.2,0.3)(-0.4,0.2)(0.3,-0.5)(-0.2,-0.3)(0.55,0.4)};
          \addplot[only marks,mark=x,mark size=2.6pt,curose,thick,forget plot]
            coordinates {(1.5,0.9)(-1.4,1.0)(1.2,-1.3)(-1.1,-1.2)(0.1,1.7)(-1.7,-0.2)};
          \addplot[cublue,very thick,domain=0:360,samples=90]
            ({cos(x)},{sin(x)});
          \node[font=\tiny,cublue,anchor=north] at (axis cs:0,-1.45) {$x_1^2+x_2^2=1$};
        \end{axis}
      \end{tikzpicture}
    \end{column}
    \begin{column}{0.5\textwidth}
      \footnotesize
      Nothing forces $z$ to be linear in the \emph{raw} inputs. Engineer polynomial features,
      exactly as in Week 2, and feed those to the same model:
      \[ z = w_1x_1^2 + w_2x_2^2 + b, \quad \vw=[1,1],\ b=-1 \]
      \[ \Rightarrow\ \text{boundary } x_1^2+x_2^2 = 1, \text{ a circle.} \]
      Inside the circle $z<0$ so $\hat y = 0$; outside, $\hat y = 1$.
      \vspace{1mm}
      \begin{keyidea}
        \footnotesize
        Higher-order terms --- $x_1^2x_2$, $x_1^3$, $x_1x_2^2$ --- give boundaries of almost
        any shape. The model stays logistic regression; only the feature vector changed.
      \end{keyidea}
      \vspace{1mm}
      With no engineered features the boundary is \emph{always} a straight hyperplane, however
      much data you have.
    \end{column}
  \end{columns}
\end{frame}

% ===============================================================
\section{The Cost Function}

\begin{frame}{Squared error will not do}
  \begin{columns}[T,onlytextwidth]
    \begin{column}{0.5\textwidth}
      \centering {\scriptsize Squared error with a sigmoid inside: two basins}\par
      \begin{tikzpicture}
        \begin{axis}[width=6.2cm,height=4.0cm,xlabel={$w$},ylabel={$J(w)$},
          xmin=-3,xmax=3,ymin=0.09,ymax=0.32,
          tick label style={font=\scriptsize},label style={font=\scriptsize},
          grid=major,grid style={culight}]
          \addplot[curose,very thick,domain=-3:3,samples=180]
            {((1/(1+exp(x)))^2 + (1/(1+exp(-x))-1)^2 + (1/(1+exp(-4*x)))^2)/6};
          \addplot[only marks,mark=*,mark size=1.6pt,cublue,forget plot]
            coordinates {(-0.19,0.1168)};
          \node[font=\tiny,cublue,anchor=east] at (axis cs:-0.45,0.122) {global};
          \node[font=\tiny,cugray,anchor=south east] at (axis cs:2.95,0.168) {local};
        \end{axis}
      \end{tikzpicture}
    \end{column}
    \begin{column}{0.48\textwidth}
      \centering {\scriptsize The logistic cost on the same data: one bowl}\par
      \begin{tikzpicture}
        \begin{axis}[width=6.2cm,height=4.0cm,xlabel={$w$},ylabel={$J(w)$},
          xmin=-3,xmax=3,ymin=0,ymax=4.3,
          tick label style={font=\scriptsize},label style={font=\scriptsize},
          grid=major,grid style={culight}]
          \addplot[cuteal,very thick,domain=-3:3,samples=180]
            {(2*ln(1+exp(-x)) + ln(1+exp(4*x)))/3};
          \addplot[only marks,mark=*,mark size=1.6pt,cublue,forget plot]
            coordinates {(-0.24,0.6549)};
        \end{axis}
      \end{tikzpicture}
    \end{column}
  \end{columns}
  \vspace{0.5mm}
  \footnotesize
  Both curves are the real cost for the three-point data set $x=(-1,1,4)$, $y=(0,1,0)$, with
  $b=0$. On the left, gradient descent started anywhere right of $w\approx0.8$ walks away from
  the better solution and settles at $J=0.1667$ instead of $0.1168$.
  \begin{keyidea}
    Squared error is convex in $\vw$ for a \emph{linear} model. Wrapping a sigmoid around it
    destroys that: $J$ acquires local minima and long flat regions, and gradient descent has
    no way to escape them. We need a different cost.
  \end{keyidea}
\end{frame}

\begin{frame}{The logistic loss}
  \footnotesize
  Write the cost as an average of a per-example \textbf{loss}, so we can design the loss on
  its own:
  \[
    J(\vw,b) = \frac{1}{m}\sum_{i=1}^{m} L\bigl(f_{\vw,b}(\xii{i}),\ \yii{i}\bigr),
    \qquad
    L(f,y) = \begin{cases}
      -\log(f) & \text{if } y = 1\\[2pt]
      -\log(1-f) & \text{if } y = 0
    \end{cases}
  \]
  \begin{columns}[T,onlytextwidth]
    \begin{column}{0.5\textwidth}
      \centering
      \begin{tikzpicture}
        \begin{axis}[width=6.2cm,height=3.9cm,xlabel={$f_{\vw,b}(\vx)$},ylabel={$L$},
          xmin=0,xmax=1,ymin=0,ymax=4.2,grid=major,grid style={culight},
          tick label style={font=\scriptsize},label style={font=\scriptsize},
          legend style={font=\tiny,at={(0.5,0.98)},anchor=north,draw=none,fill=white}]
          \addplot[cuteal,very thick,domain=0.01:0.995,samples=120]{-ln(x)};
          \addlegendentry{$y=1$: $-\log f$}
          \addplot[curose,very thick,domain=0.005:0.99,samples=120]{-ln(1-x)};
          \addlegendentry{$y=0$: $-\log(1-f)$}
        \end{axis}
      \end{tikzpicture}
    \end{column}
    \begin{column}{0.46\textwidth}
      \begin{center}\scriptsize
      \renewcommand{\arraystretch}{1.15}
      \begin{tabular}{@{}rrr@{}}
        \toprule
        $f$ & $L$ if $y=1$ & $L$ if $y=0$\\
        \midrule
        $0.99$ & $0.0101$ & $4.6052$\\
        $0.90$ & $0.1054$ & $2.3026$\\
        $0.50$ & $0.6931$ & $0.6931$\\
        $0.10$ & $2.3026$ & $0.1054$\\
        \bottomrule
      \end{tabular}
      \end{center}
      \vspace{1mm}
      A confidently wrong prediction costs without bound; a confidently right one costs almost
      nothing. That unbounded penalty is what pushes $f$ towards the correct end, and what
      makes $J$ convex.
    \end{column}
  \end{columns}
\end{frame}

\begin{frame}{The simplified cost function}
  \small
  The two branches can be written as one expression, because $y$ is always $0$ or $1$:
  \[
    \boxed{\;L(f,y) \;=\; -\,y\log(f) \;-\; (1-y)\log(1-f)\;}
  \]
  \footnotesize
  Averaging over the training set gives the cost we will minimise, writing
  $f^{(i)}=f_{\vw,b}(\xii{i})$:
  \[
    J(\vw,b) = -\frac{1}{m}\sum_{i=1}^{m}
      \Bigl[ \yii{i}\log\bigl(f^{(i)}\bigr) + \bigl(1-\yii{i}\bigr)\log\bigl(1-f^{(i)}\bigr) \Bigr]
  \]
  \begin{columns}[T,onlytextwidth]
    \begin{column}{0.48\textwidth}
      \textbf{Check both cases.}
      \begin{itemize}
        \item $y=1$: the second term has factor $(1-1)=0$, leaving $-\log(f)$.
        \item $y=0$: the first term has factor $0$, leaving $-\log(1-f)$.
      \end{itemize}
      One of the two terms is always switched off, so nothing has changed except that the
      formula is now a single differentiable expression. Note there is no $\tfrac12$: it only
      ever cancelled the $2$ from differentiating a square.
    \end{column}
    \begin{column}{0.48\textwidth}
      \begin{keyidea}
        \footnotesize
        This is not an arbitrary choice. It is the negative log-likelihood of the data under
        the model, so minimising $J$ is \textbf{maximum likelihood estimation}. It is also
        called \emph{binary cross-entropy}, and it is provably convex in $(\vw,b)$.
      \end{keyidea}
    \end{column}
  \end{columns}
\end{frame}

% ===============================================================
\section{Gradient Descent for Logistic Regression}

\begin{frame}{The update rule --- familiar, but not the same}
  \begin{defnbox}[Gradient descent for logistic regression]
    Repeat until convergence, updating all $n+1$ parameters \emph{simultaneously}:
    \[
      \begin{aligned}
        w_j &:= w_j - \alpha\,\frac{1}{m}\sum_{i=1}^{m}
          \bigl(f_{\vw,b}(\xii{i})-\yii{i}\bigr)\xjj{i}{j}, \qquad j=1,\dots,n\\[2pt]
        b &:= b - \alpha\,\frac{1}{m}\sum_{i=1}^{m}\bigl(f_{\vw,b}(\xii{i})-\yii{i}\bigr)
      \end{aligned}
    \]
  \end{defnbox}
  \footnotesize
  These are algebraically identical to Week 2 --- the $\log$ and the sigmoid conspire so that
  the derivative comes out this clean. But
  \[
    \text{linear: } f_{\vw,b}(\vx)=\vw\cdot\vx+b
    \qquad\text{versus}\qquad
    \text{logistic: } f_{\vw,b}(\vx)=\frac{1}{1+e^{-(\vw\cdot\vx+b)}} .
  \]
  \begin{pitfall}
    Same formula, different algorithm. Everything from Week 2 still applies: simultaneous
    updates, feature scaling, monitoring $J$ against iteration, the $\times3$ search for
    $\alpha$ --- and a vectorized implementation.
  \end{pitfall}
\end{frame}

\begin{frame}{Four iterations, worked by hand}
  \footnotesize
  Data: $x = (1,2,3,4)$, $y = (0,0,1,1)$; one feature, $m=4$, $\alpha=0.5$, starting from
  $w=0$, $b=0$.
  \begin{center}\scriptsize
  \renewcommand{\arraystretch}{1.2}
  \begin{tabular}{@{}crrrrr@{}}
    \toprule
    iter & $\partial J/\partial w$ & $\partial J/\partial b$ & $w$ & $b$ & $J(w,b)$\\
    \midrule
    $0$ & --- & --- & $0.0000$ & $0.0000$ & $0.6931$\\
    $1$ & $-0.5000$ & $+0.0000$ & $0.2500$ & $0.0000$ & $0.6250$\\
    $2$ & $-0.0578$ & $+0.1487$ & $0.2789$ & $-0.0744$ & $0.6124$\\
    $3$ & $-0.0525$ & $+0.1476$ & $0.3051$ & $-0.1482$ & $0.6003$\\
    $4$ & $-0.0512$ & $+0.1452$ & $0.3307$ & $-0.2207$ & $0.5885$\\
    \bottomrule
  \end{tabular}
  \end{center}
  \vspace{1mm}
  \begin{columns}[T,onlytextwidth]
    \begin{column}{0.5\textwidth}
      \begin{itemize}
        \item At $w=b=0$ every $f=g(0)=0.5$, so $J=-\log 0.5 = 0.6931$ --- the cost of
              guessing.
        \item $J$ falls on every iteration. That is the check from Week 2, unchanged.
        \item Progress is slow: after four steps the boundary $-b/w$ sits at $0.67$, far from
              the sensible $2.5$. Hundreds of iterations are normal.
      \end{itemize}
    \end{column}
    \begin{column}{0.46\textwidth}
      \begin{keyidea}
        \footnotesize
        $J = 0.6931 = \log 2$ is the cost of a model that predicts $0.5$ for everything. Use
        it as a baseline: if your trained $J$ is not clearly below $\log 2$, the model has
        learned nothing.
      \end{keyidea}
    \end{column}
  \end{columns}
\end{frame}

\begin{frame}{In code}
  \footnotesize
  \begin{block}{The whole model in four functions}
    \texttt{def sigmoid(z):\ \ \ \ return 1 / (1 + np.exp(-z))}\\[2pt]
    \texttt{def predict(X, w, b):\ return sigmoid(X @ w + b)\ \ \ \ \ \ \ \ \# shape (m,)}\\[2pt]
    \texttt{def cost(X, y, w, b):}\\
    \texttt{\ \ \ \ f = predict(X, w, b)}\\
    \texttt{\ \ \ \ return -np.mean(y*np.log(f) + (1-y)*np.log(1-f))}\\[2pt]
    \texttt{def gradient(X, y, w, b):}\\
    \texttt{\ \ \ \ err = predict(X, w, b) - y}\\
    \texttt{\ \ \ \ return X.T @ err / len(y),\ \ err.mean()}
  \end{block}
  \begin{columns}[T,onlytextwidth]
    \begin{column}{0.5\textwidth}
      \begin{pitfall}[\texttt{log(0)} $= -\infty$]
        With large $|z|$, \texttt{sigmoid} rounds to exactly $0$ or $1$ and the cost becomes
        \texttt{nan}. Production code uses a numerically stable form; for the lab, clip $f$
        into $[10^{-15},\,1-10^{-15}]$.
      \end{pitfall}
    \end{column}
    \begin{column}{0.46\textwidth}
      \small
      Compare with Week 2: only \texttt{predict} and \texttt{cost} changed. \texttt{gradient}
      is character-for-character the linear-regression version, because the mathematics says
      it must be.
    \end{column}
  \end{columns}
\end{frame}

% ===============================================================
\section{Wrap-up and Assessment}

\begin{frame}{Summary of Part 1}
  \footnotesize
  \begin{columns}[T,onlytextwidth]
    \begin{column}{0.5\textwidth}
      \begin{block}{Concepts}
        \begin{itemize}
          \item Classification predicts a category; $y\in\{0,1\}$.
          \item Linear regression fails: unbounded output, and the boundary moves when a
                far-away correct point is added.
          \item $f_{\vw,b}(\vx)=g(\vw\cdot\vx+b)$ estimates $P(y=1\mid\vx)$.
          \item Thresholding at $0.5$ makes the boundary $\vw\cdot\vx+b=0$; polynomial
                features bend it.
          \item Squared error is non-convex here; logistic loss is convex.
          \item The gradient formula is unchanged from Week 2 --- $f$ is not.
        \end{itemize}
      \end{block}
    \end{column}
    \begin{column}{0.46\textwidth}
      \begin{block}{Formulas}
        \[ g(z)=\frac{1}{1+e^{-z}} \]
        \[ J = -\frac1m\sum_i\bigl[\yii{i}\log f^{(i)} + (1-\yii{i})\log(1-f^{(i)})\bigr] \]
        \[ \frac{\partial J}{\partial w_j}=\frac1m\sum_i\bigl(f^{(i)}-\yii{i}\bigr)\xjj{i}{j} \]
      \end{block}
      \begin{block}{Next}
        \footnotesize Part 2: what goes wrong when a model fits the training data too well,
        and how regularization fixes it --- for both regression and classification.
      \end{block}
    \end{column}
  \end{columns}
\end{frame}

\begin{frame}{Descriptive questions}
  \footnotesize
  \begin{enumerate}
    \item Define binary classification, and state which class is conventionally labelled $1$.
          Give two examples from outside the slides. \hfill\clo{4}
    \item Using the tumour example, explain two distinct reasons linear regression is
          unsuitable for classification. \hfill\clo{6}
    \item State the sigmoid function, its range, and its value at $z=0$. Prove that
          $g(-z)=1-g(z)$. \hfill\clo{5}
    \item Write the logistic regression model and state precisely what its output means. If
          $f_{\vw,b}(\vx)=0.3$, what does the model believe about this example? \hfill\clo{5}
    \item Show that thresholding $f_{\vw,b}(\vx)$ at $0.5$ is equivalent to thresholding
          $\vw\cdot\vx+b$ at $0$. \hfill\clo{5}
    \item For $\vw=[1,1]$ and $b=-3$, give the equation of the decision boundary and compute
          $f_{\vw,b}(\vx)$ at $\vx=(1,1)$ and $\vx=(2,2)$. \hfill\clo{6}
  \end{enumerate}
\end{frame}

\begin{frame}{Descriptive questions (continued)}
  \footnotesize
  \begin{enumerate}\setcounter{enumi}{6}
    \item Explain how polynomial features let logistic regression produce a circular decision
          boundary. Give $\vw$ and $b$ for the boundary $x_1^2+x_2^2=4$. \hfill\clo{6}
    \item Why is squared error unsuitable as the cost for logistic regression? Refer to
          convexity and to the behaviour of gradient descent. \hfill\clo{6}
    \item Write the logistic loss in its two-case form and in its single-line form, and show
          the two agree for $y=0$ and for $y=1$. \hfill\clo{6}
    \item Compute $L(f,y)$ for $(f,y)=(0.9,1)$, $(0.1,1)$ and $(0.1,0)$. Comment on the
          relative sizes. \hfill\clo{5}
    \item Write the gradient descent updates for logistic regression. State what is identical
          to linear regression and what is not. \hfill\clo{6}
    \item An untrained logistic model has $J=\log 2 \approx 0.693$. Explain why, and how you
          would use this number when debugging a training run. \hfill\clo{5}
  \end{enumerate}
\end{frame}

\begin{frame}{Design questions}
  \footnotesize
  \begin{enumerate}
    \item \textbf{Design the threshold.} A Chattogram hospital will use your model to decide
          which patients get a biopsy. Missing a malignant tumour is far worse than an
          unnecessary biopsy. Choose a threshold, justify it, and describe how you would
          present the trade-off to the clinicians who must sign off on it. \hfill\clo{12}
    \item \textbf{Design a feature set for a decision boundary.} Sketch two-dimensional data
          for which no straight line separates the classes, then specify the engineered
          features that would let logistic regression separate it. State the price you pay
          for adding them. \hfill\clo{12}
    \item \textbf{Design a numerical-safety layer.} Your \texttt{cost} function returns
          \texttt{nan} after 300 iterations. Design a defensive implementation: what you clip,
          what you assert, and what you log so the cause is identifiable from the log alone.
          \hfill\clo{10}
    \item \textbf{Design a test suite.} Propose four tests for a logistic regression
          implementation that would catch, respectively: a missing sigmoid, a sign error in
          the gradient, non-simultaneous updates, and unscaled features. For each, state the
          input and the observable symptom. \hfill\clo{12}
    \item \textbf{Design for accountability.} A bank must tell a rejected applicant why. Given
          a trained logistic model, describe what you can honestly report from $\vw$, $b$ and
          $\vx$, and one thing that would be dishonest to claim. \hfill\clo{10}
  \end{enumerate}
\end{frame}

\begin{frame}{Reading}
  \begin{block}{Primary text}
    G.~James, D.~Witten, T.~Hastie and R.~Tibshirani,
    \emph{An Introduction to Statistical Learning}, 2nd ed., Springer, 2021.\\[2pt]
    \hfill{\small$\rightarrow$ \textbf{\S4.1--4.3} --- why not linear regression, the logistic
    model, the logistic function, estimating the coefficients by maximum likelihood, and
    multiple logistic regression.}
  \end{block}
  \begin{block}{Supplementary}
    C.~M.~Bishop, \emph{Pattern Recognition and Machine Learning}, Springer, 2006.\\[2pt]
    \hfill{\small$\rightarrow$ \textbf{\S4.3.2} --- logistic regression, the cross-entropy
    error function, and the derivation of the gradient that makes it look like the linear
    case.}
  \end{block}
  \vspace{2mm}
  \begin{center}
    \small\color{cugray}
    Lab (CSE 816, Week 3): classification; the sigmoid function; decision boundary;
    logistic loss and cost.\\
    Practice quizzes: \emph{Classification with logistic regression}; \emph{Cost function for
    logistic regression}.
  \end{center}
\end{frame}

\end{document}
